\section{Conductors and finite multiplication-norm transport}\label{sec:norms}
\noindent\textit{Source anchor T: foundation v0.1--v0.4. These transport results concern analytic elliptic functions or their divisor classes, not an arithmetic evaluation of their coefficients.}

\subsection{Exactly normalized elliptic norms}
On $E_t=\mathbb C/\Lambda_t$ define $N_nf(w)=\prod_{nz=w}f(z)$ for integers $n\ge1$, with all $n^2$ preimages counted. The product is independent of the representative lift. Root-of-unity averaging and vertical residue splitting in \eqref{eq:theta-product} give
\begin{equation}\label{eq:qnorm}
 N_nq_t(w)=\begin{cases}q_t(w),&5\nmid n,\\1,&5\mid n.\end{cases}
\end{equation}
In the prime-to-five case the two product identities are
\[
 \prod_{a=0}^{n-1}\theta_0(z+a/n;t)=\theta_0(nz;nt),\qquad
 \prod_{b=0}^{n-1}\theta_0(nz+bt;nt)=\theta_0(nz;t).
\]
Multiplication by five permutes the real residues in the smoothing denominator. When $5\mid n$, real averaging retains only frequencies already deleted by fifth smoothing. This proves the scalar one in both cases, rather than a norm formula up to a root of unity.

The point $S=s\bmod\Lambda_t$ has infinite order by rational independence. Put
\[
 \psi(w)=\frac{q_t(w)^2}{q_t(2w)q_t(2w+s)}.
\]
\begin{proposition}[All finite norms of the return kernel]\label{prop:allnorms}
For every integer $n\ge1$,
\[
 N_n\psi(w)=
 \begin{cases}
 1,&5\mid n,\\[2pt]
 \displaystyle\frac{q_t(w)^2}{q_t(2w)q_t(2w+ns)},&5\nmid n,\ n\text{ odd},\\[7pt]
 \displaystyle\frac1{q_t(w)^2q_t(w+ks)^4},&5\nmid n,\ n=2k.
 \end{cases}
\]
\end{proposition}
\begin{proof}
For odd $n$, doubling permutes $E_t[n]$. For $n=2k$, it maps $E_t[n]$ to $E_t[k]$ with four lifts per point, so each denominator norm is a fourth power. Apply \eqref{eq:qnorm} and cancel the numerator. If $5\mid n$ then all these norms are one, including the even case since $5\mid k$.
\end{proof}
For $n=2k+1$ prime to five, $B_k(w)=\prod_{j=0}^{k-1}q_t(2w+(2j+1)s)$ gives
\[
 N_{2k+1}\psi(w)=\psi(w)\frac{B_k(w)}{B_k(w+s)}.
\]
For $n=2k$ prime to five, $C_k(w)=\prod_{j=0}^{k-1}q_t(w+js)^4$ gives
\[
 N_{2k}\psi(w)=q_t(w)^{-6}\frac{C_k(w)}{C_k(w+s)}.
\]
Finite products therefore retain explicit endpoint multipliers. Closing a permutation of labels does not remove them or fix the scalar normalization of a separate orbit.

\subsection{The exact information lost by doubling}
Quotient nonzero meromorphic elliptic functions by constants and translation coboundaries $A(w+s)/A(w)$. Set $e=[q_t]$, $p=[\psi]$ in additive notation. These generate a free $\mathbb Z^2$: the divisor charges on the disjoint infinite translation orbits of $0$ and $-S/2$ are respectively $(4,4)$ and $(0,-4)$. A finite divisor coboundary has charge zero on every orbit by telescoping, so a relation $ae+bp=0$ forces $b=0$ then $a=0$.

The norm matrices on this free subgroup are
\[
 N_n=I\quad(n\text{ odd},5\nmid n),\qquad
 N_n=\begin{pmatrix}1&-6\\0&0\end{pmatrix}\quad(n\text{ even},5\nmid n),\qquad
 N_n=0\quad(5\mid n).
\]
The middle matrix is a split integral projector with image $\mathbb Z(1,0)$ and kernel $\mathbb Z(6,1)$. Its primitive kernel class $6e+p$ is nonzero, as its two charges are $28,-4$. Thus repeated full doubling cannot recover the deleted component. The residue-ten matrix law is multiplicative; its finite coefficient statement belongs to the inherited formal core, whereas the application of divisors to analytic functions is an ordinary argument.

\subsection{The fifth-isogeny differential retains its constant}
Let $a=\wp_t(1/5)$, $b=\wp_t(2/5)$ and $L_t=(q_t'/q_t)'$. The enlarged lattice is $\Lambda_t^+=\tfrac15\mathbb Z+\mathbb Zt$, with $\wp_+(z)=25\wp_{5t}(5z)$. The degree-five quotient satisfies
\begin{equation}\label{eq:isogeny-differential}
 L_t(z)=\wp_+(z)-5\wp_t(z)+2(a+b).
\end{equation}
At zero $L_t$ has principal part $-4/z^2$; at the four nonzero kernel points it has $+1/(z-P_j)^2$, with no simple-pole terms. Thus $L_t-\wp_++5\wp_t$ is a constant elliptic function. The canonical formula gives $q_t(z)=(a-b)^2z^4[1+(a+b)z^2+O(z^4)]$, fixing that constant as $2(a+b)$. The degree-five quotient is equivalently the paired sum
\[
 \wp_+(z)=\wp_t(z)+\sum_{j=1}^{4}[\wp_t(z+j/5)-\wp_t(j/5)].
\]
The analytic identity determines a rational isogeny map once elliptic coordinates are chosen. It does not assert that the period-dependent parameter is an algebraic number or identify this cubic-period curve with the separate equianharmonic CM curve.
